\chapter{Tests and Evaluation}
\label{chapter:evaluation}

Throughout this project Smache has been implemented in Python and integrated into the Peergrade.io-platform. Instructions for where to find and use the library is available in appendix~\ref{chapter:source-code}. The library is also open sourced on GitHub available on the following URL:

\url{https://github.com/anderslime/smache}

Besides the Python implementation of Smache, the source code also includes automated tests with a test coverage of 97\% and a small benchmark framework to execute and plot the results of performance tests. The automated tests are automatically verified by a continuous integration application, Travis, that runs the full test suite when new versions of the code are merged into the source code. The results of the automated tests are available at \url{https://travis-ci.org/anderslime/smache}.

The rest of this chapter will lay out results of experiments made to test the assumptions about the solution. The experiments are performed on a MacBook Pro (Retina, 15-inch, Early 2013) with an 2.4 GHz Intel Core i7 (4 cores) and 8 GB 1600 MHz DDR3 RAM. The evaluation seeks to answer the following questions:

\begin{itemize}
  \item What changes are required to modify existing web application to use Smache? (section~\ref{sec:changes-required-to-cache-with-smache})
  \item What is the performance impact for procedures updating data in the primary storage, when cachable functions are introduces with Smache? (section~\ref{sec:performance-impact-of-existing-procedures})
  \item How efficient are Smache updating cached objects? (section~\ref{sec:update-throughput})
\end{itemize}

The chapter will end with an evaluation of the system based on the test results and the initial requirements~\ref{sec:evaluation}.

\section{Changes Required to use Smache}
\label{sec:changes-required-to-cache-with-smache}

Smache was integrated into the Peergrade.io platform, which is an existing web application with a single code base deployed to multiple web servers. To test the adaptability we applied the library to cached functions of different granularities such as functions returning a single value, functions returning a data structure and functions returning the HTML served to the user. From this integrations we experienced an easy integration, but some of the functions had to be split up to to cache optimally .
For example some of the functions had statistical calculations that iterated using \verb$for$-loops and made time consuming calculations in each iteration as seen in code snippet~\ref{code:running-example-non-cachable}, which were not able to be cached individually. To make this function optimally cachable, we would have to extract it into a method such as in the running example in code snippet~\ref{code:running-example}. The running example are then made cachable by adding cachable declaration using Python-decorators after which it looks as in code snippet~\ref{code:basic-caching-running-example}.
The HTTP-responses were not directly cachable, since they work on request-objects specific to the client making the request. In order to cache the responses we extracted the code that generated the HTML-fragment served as the HTTP-body.

In order to use Smache, the system also require additional components added to the architecture. To achieve the guarantees described in this thesis, the system must add the following components to the architecture, also illustrated on figure~\ref{fig:smache-architecture-changes}:

\begin{itemize}
  \item Cache database to store the cached values and metadata
  \item Background Workers with a Queue-system to have asynchronous invalidation and updates
\end{itemize}

\begin{figure*}[ht!]
  \centering
  \includegraphics[width=1.0\linewidth]{figures/smache-architecture-changes.pdf}
  \caption{Additional architectural components required to use Smache compared to a normal web application including clients, web application servers and a primary storage.}
  \label{fig:smache-architecture-changes}
\end{figure*}

Even though this is optimal, Smache has been implemented such that it is easy leave them out. The cache database can for example be replaced by the primary storage that is already a part of the system. The background workers can also be left out and replaced by synchronous invalidation and updates that are executed in the same process as the web servers, but this would affect the performance of the web applications as seen in the results in section~\ref{sec:performance-impact-of-existing-procedures}.

\begin{code}{Code where Smache can cache the course score, but not the individual participant scores.}
  \input{code/running_example_non_cachable.py}
  \label{code:running-example-non-cachable}
\end{code}

% section changes-required-to-cache-with-smache end

\section{Performance Impact of Existing Operations}
\label{sec:performance-impact-of-existing-procedures}

One of the requirements of the system is that the performance of existing operations should not be affected negatively by introducing the caching system. On the requests that changes underlying data, the application will notify and execute invalidation in the same process, which could affect the response time of the update request. The performance of invalidation have been evaluated using a test that uses both \emph{direct dependencies} and \emph{lazy dependencies}.

The results of this test, visualized on figure~\ref{fig:graph-update-requests}, indicates how a synchronous and asynchronous write-through mechanism impacts the time taken to update underlying data. If we consider the asynchronous mechanism used by Smache, it introduces a small constant overhead of around $2$ ms as seen on the exact results shown on table~\ref{tab:update-requests}. The synchronous mechanism introduces an overhead that increases with the execution time of the cached function, which can become critical for the user experience with long running computations.

\begin{figure*}[ht!]
  \centering
  \includegraphics[width=1.0\linewidth]{figures/results/graph_update_requests.pdf}
  \caption{Impact of introducing synchronous and asynchronous write-through in a case, where the update affects a single cached object instance with variable execution time.}
  \label{fig:graph-update-requests}
\end{figure*}

\begin{table}[ht!]
  \centering
  \begin{tabular}{llll}
    \hline
    Function duration & Without caching & Smache (async.) & Synchronous \\
    \hline
    50 ms  & 6.2 ms (1.0x)  & 7.9 ms (1.27x)  &  143.2 ms (23.0x) \\
    100 ms & 6.3 ms (1.0x)  & 7.9 ms (1.25x)  &  146.9 ms (23.4x) \\
    200 ms & 6.4 ms (1.0x)  & 8.1 ms (1.26x)  &  168.2 ms (26.4x) \\
    400 ms & 6.3 ms (1.0x)  & 8.0 ms (1.25x)  &  207.1 ms (31.2x) \\
    \hline
  \end{tabular}
  \caption{The time taken updating underlying data that affects a single cached object for an application using Smache compared to a system using synchronous write-through invalidation.}
  \label{tab:update-requests}
\end{table}

% section performance-impact-of-existing-procedures end

\section{Update Throughput}
\label{sec:update-throughput}

To evaluate the performance and scalability of the concurrency used in the data update propagation, we will measure the update throughput while increasing the number of workers, where the number of workers is the amount updates the system is able to process at the same time. While it would be possible to measure against examples in the context of the Peergrade.io, we have chosen to measure against patterns of cached functions to make the experiments easy to reproduce for future research. The tests will be performed against a use case, where many cached functions depends on a single underlying data resource (section~\ref{subsec:many-cached-object-instances} and afterwards a case where multiple cached functions depends on each other in a nested hierarchy (section~\ref{subsec:nested-cached-functions}).

\subsection{Many Cached Object Instances}
\label{subsec:many-cached-object-instances}

In this use case we will test a use case, where multiple instances of the same cached function depends on the same underlying data such that when the underlying data updates then all the cached objects instances for the cached function must be updated. The use case is illustrated on figure~\ref{fig:test-multi-deps}.

\begin{figure*}[ht!]
  \centering
  \includegraphics[width=1.0\linewidth]{figures/test-multi-deps.pdf}
  \caption{Illustration of the dependencies of the cached function for the "Many Cached Functions" test case}
  \label{fig:test-multi-deps}
\end{figure*}

To simulate slow functions the tests use the \verb$sleep$-method from Python's \verb$time$-module such that each function uses $400$ ms to execute.

The results of the tests are visualized figure~\ref{fig:graph_throughput_workers_multi} by a graph showing the number of updates the system can process per second (throughput) for different amount of workers. This graph shows how concurrency affects the throughput of the cache system, where we see a significant increase in throughput every time the number of workers are increased. We also see that the impact of adding workers are becomes higher when there are more cached objects even though it also has an impact for small amounts (around $10$ cached objects).

\begin{figure*}[ht!]
  \centering
  \includegraphics[width=1.0\linewidth]{figures/results/graph_throughput_workers_multi.pdf}
  \caption{How the system scales with many cached functions depending on the same underlying data while the number of workers is increased}
  \label{fig:graph_throughput_workers_multi}
\end{figure*}

% subsection many-cached-object-instances end

\subsection{Nested Cached Functions}
\label{subsec:nested-cached-functions}

To evaluate the performance on nested cached functions, we apply the tests on a setup, where multiple cached functions depend on each other as illustrated on figure~\ref{fig:test-nested-deps}.

\begin{figure*}[ht!]
  \centering
  \includegraphics[width=0.1\linewidth]{figures/test-nested-deps.pdf}
  \caption{Illustration of the dependencies of the cached function for the "Nested Cached Functions" test case}
  \label{fig:test-nested-deps}
\end{figure*}

Compared to the impact on the use case with many cached functions (section~\ref{subsec:many-cached-object-instances}), the results on figure~\ref{fig:graph_throughput_workers_nested} shows that the impact of concurrent updates are not as high and that the impact highly depends on the depth of the cached functions.

\begin{figure*}[ht!]
  \centering
  \includegraphics[width=1.0\linewidth]{figures/results/graph_throughput_workers_nested.pdf}
  \caption{How the system scales with nested cached functions while the number of workers is increased}
  \label{fig:graph_throughput_workers_nested}
\end{figure*}

Another reason why the use case with nested cached functions are interesting is because it cannot be parallelized by scheduling algorithms that assumes a given cached function is only computed once at the same time. These algorithms would process the cached functions in this test case sequentially, which corresponds to the result of using 1 worker. If we compare the throughput of using multiple workers to the throughput of using a single worker as seen on table~\ref{tab:nested-relative-results}, we see that there is still a gain in throughput from processing multiple functions at once.

\begin{table}[ht!]
  \centering
  \begin{tabular}{llll}
    \hline
    Cached functions & Sequentially & Concurrently & Improvement \\
    \hline
    2 & 0.80 updates/sec & 1.57 updates/sec & 1.97x \\
    4 & 0.94 updates/sec & 1.88 updates/sec & 2.01x \\
    8 & 0.96 updates/sec & 2.43 updates/sec & 2.55x \\
    \hline
  \end{tabular}
  \caption{Improvement of processing nested cached functions concurrently (using 8 workers) compared to sequentially (using 1 worker).}
  \label{tab:nested-relative-results}
\end{table}

\section{Evaluation}
\label{sec:evaluation}

The final solution, Smache, solves problem described in section~\ref{sec:problem} by introducing a caching system that provides automatic invalidation and on-write updates to ensure that the cached fragments are always served instantly. Smache allows the programmer to define whether a given cached function should serve stale values or recompute a value if it is not fresh. To relate to existing solutions table~\ref{table:smache-evaluation} shows how Smache relates to the cache evaluation criteria (see appendix~\ref{chapter:comparison-of-caching-approaches-including-smache} for the table including existing approaches).

\begin{table}[ht!]
  \scriptsize
  \doublespacing
  \centering
  \hspace*{-1.5cm}
  \begin{tabular}{lcccccc}
{} & {
  \textbf{Consistency}
} & {
  \twolinecell{1.2cm}{Strict}{Freshness}
} & {
  \twolinecell{1.2cm}{In-place}{Updates}
} & {
  \threelinecell{1.4cm}{Always}{Immediate}{Response}
} & {
  \twolinecell{1.4cm}{No Cache}{Management}
} & {
  \textbf{Adaptability}
} \\
  \hline
  \textbf{Arbitrary Content} & & & & & & \\
  \multicolumn{7}{c}{...} \\
  Smache               & \no & \no  & \yes & \yes & \opt\sss{*} & \med \\[7pt]
  Smache w/o staleness & \no & \yes & \yes & \no  & \opt\sss{*} & \med \\[7pt]
  \multicolumn{7}{c}{...} \\
  \hline
  \multicolumn{7}{l}{*) Dependencies to underlying data must be declared}
  \end{tabular}
  \caption{How Smache relates to the cache evaluation criteria}
  \label{table:smache-evaluation}
\end{table}

In order to make the Smache library adaptable for existing application it has been designed and implemented to be flexible and introduce as few new components as possible. The architectural components are made as extensions such that they easily can be changed or left out if some of the guarantees are not needed in the given context. The interface for making functions cachable is not fully transparent, but it relies on the developer declaring dependencies. Smache could be extended with a transparent interface, where the dependencies are automatically detected, but this would also affect the flexibility, the control of the programmer and make it more difficult to understand when a given cached functions are invalidated. We therefore suggest to use a tool that helps the programmer identify dependencies instead of removing the control or a solution where the programmer optionally is able to declare dependencies. This also relates to the requirement of making the system maintainable and understandable for the developer.

The test results of this chapter indicates that the Smache solution does not introduce performance regressions to the system. The requests that updates underlying data are normally affected when the application uses a basic write-through invalidation, but Smache only introduces a small constant overhead of a few milliseconds ($< 3 ms$) by invalidating and updating using an asynchronous technique that uses timestamp invalidation for correctness. The update propagation algorithm allows concurrent updates that in some cases results in a higher throughput than scheduling algorithms. The algorithm uses the timestamp invalidation technique to ensure correctness such that the updates can be executed concurrently, which means we achieve a simpler system that does not require scheduling algorithms to achieve a high throughput and we are able to scale the throughput by adding more workers, but it also means that algorithm can risk executing the same cached functions more than once and thereby using more CPU power than necessary. We can therefore say that the system is efficient with relation to the performance of the web application and the throughput, but not with relation to CPU usage.

Smache does not control the space usage, which means the presence of the cached values depends on how the cache database controls space. The optimal choice for a cache database is a in-memory database, but since memory is expensive, it could be necessary to change to a disk-based database if the space used becomes too expensive.

The system has been designed with fault-tolerant measures such as isolation and retries of invalidation and updates, but to keep the solution simple, the system does not apply measures to ensure the availability of the cache system. More specifically the design does not implement measures to prevent or recover from faults happening in the queuing system or the cache database, which means they are seen as a single point of failure of the system. To run the application correctly it must be ensure that the queue and cache database are highly available.

% section evaluation end
